%------------------------------------------------------------------------------%

\usepackage{apresentacao}
\usepackage{mathconfig}
\usepackage{tabto}

\makeatletter
\graphicspath  {{./figs/}}
\def\input@path{{./figs/}}
\makeatother

%------------------------------------------------------------------------------%
\title {Volumes por Cascas Cilíndricas}
\author{\large Luis Alberto D'Afonseca}
\date  {Integrais e Séries}

%------------------------------------------------------------------------------%
\begin{document}

%------------------------------------------------------------------------------%
\begin{frame}
    \titlepage
\end{frame}

%------------------------------------------------------------------------------%
\section{Casca Cilíndricas}

%------------------------------------------------------------------------------%
\framedgraphic{Casca Cilíndrica}{casca_cilindrica}{}

%------------------------------------------------------------------------------%
\begin{frame}{Volume de uma Casca Cilíndrica}
\begin{align*}
  \onslide<+->{V & = V_2 - V_1                                          \\[2mm]}
  \onslide<+->{  & = \pi r_2^2 h - \pi r_1^2 h                          \\[2mm]}
  \onslide<+->{  & = \pi \left(r_2^2 - r_1^2\right) h                   \\[2mm]}
  \onslide<+->{  & = \pi \left(r_2 + r_1\right)\left(r_2 - r_1\right) h \\[2mm]}
  \onslide<+->{  & = 2 \pi \frac{r_2 + r_1}{2}\left(r_2 - r_1\right) h  \\[2mm]}
  \onslide<+->{  & = 2 \pi r\, \Delta r\, h}
\end{align*}
\onslide<6->{
onde fizemos
\(\; r = \frac{r_2 + r_1}{2} \;\) e
\(\; \Delta r = r_2 - r_1    \;\)}
\end{frame}

%------------------------------------------------------------------------------%
\begin{frame}{Volume de uma Casca Cilíndrica}
  \begin{align*}
      V & = (\text{circunferência})\,(\text{altura})\,(\text{espessura}) \\[3mm]
        & = (2 \pi r) \; (h) \; (\Delta x)
  \end{align*}

\end{frame}

%------------------------------------------------------------------------------%
\section{Volumes por Cascas Cilíndricas}

%------------------------------------------------------------------------------%
\framedgraphic{Sólido de Revolução}{solido_revolucao}{}
\framedgraphic{Volume Sólido de Revolução}{solido_revolucao_volume_1}{}
\framedgraphic{Volume Sólido de Revolução}{solido_revolucao_volume_2}{}
\framedgraphic{Volume Sólido de Revolução}{solido_revolucao_volume_3}{}

%------------------------------------------------------------------------------%
\begin{frame}{Volumes por Cascas Cilíndricas}
  \begin{align*}
    \onslide<+->{V & = \int_a^b dV                                                           \\[2mm]}
    \onslide<+->{  & = \int_a^b (\text{circunferência})\,(\text{altura})\,(\text{espessura}) \\[2mm]}
    \onslide<+->{  & = \int_a^b 2\pi x \; f(x) \; dx}
  \end{align*}
\end{frame}

%------------------------------------------------------------------------------%
\begin{frame}{Volumes por Cascas Cilíndricas}
  O volume do sólido $S$ obtido pela rotação em torno do eixo $y$
  da região
  \[
    R = \big\{(x, y)\colon\; 0 \leq y \leq f(x),\; a \leq x \leq b \big\}
  \]
  é calculado pela integral
  \[
    V = \int_a^b 2\pi\, x\, f(x)\, dx
  \]
\end{frame}

%------------------------------------------------------------------------------%
\section{Exemplo 1}

\begin{frame}{Exemplo 1}
  Encontre o volume do sólido obtido pela rotação em torno do eixo $y$\\
  da região delimitada por
  \[
    y = 2x^2 - x^3
    \qquad\text{e}\qquad
    y = 0
  \]
\end{frame}

\framedgraphic{Exemplo 1 -- Esboço}{Volume_cascas_exemplo_1}{}

\begin{frame}{Exemplo 1 -- Casca Cilíndrica}
  Raio \tabto{30mm}\emph{$r(x) = x$}

  \pause
  \vspace{\baselineskip}
  Circunferência \tabto{30mm}\emph{$2\pi x$}

  \pause
  \vspace{\baselineskip}
  Altura \tabto{30mm}\emph{$h(x) = 2x^2 - x^3$}
\end{frame}

\begin{frame}{Exemplo 1 -- Volume}
  \begin{align*}
    \onslide<+->{V & = \int_0^2 2\pi x\left(2x^2 - x^3\right) dx                       \\[2mm]}
    \onslide<+->{  & = 2\pi \int_0^2 2x^3 - x^4 dx                                     \\[2mm]}
    \onslide<+->{  & = 2\pi \left( 2\frac{x^4}{4}- \frac{x^5}{5} \right) \evalat{0}{2} \\[2mm]}
    \onslide<+->{  & = 2\pi \left( 2^3 -\frac{2^5}{5}\right)}
    \onslide<+->{  \;=\; 2\pi \left( \frac{5\times 8}{5} -\frac{32}{5}\right)}
    \onslide<+->{  \;=\; 2\pi \frac{8}{5}}
    \onslide<+->{  \;=\; \frac{16\pi}{5}}
  \end{align*}
\end{frame}

%------------------------------------------------------------------------------%
\section{Exemplo 2}

\begin{frame}{Exemplo 2}
  Encontre o volume do sólido obtido pela rotação em torno do eixo $y$\\
  da região entre
  \[
    y = x
    \qquad\text{e}\qquad
    y = x^2
  \]
\end{frame}

\framedgraphic{Exemplo 2 -- Esboço}{Volume_cascas_exemplo_2}{}

\begin{frame}{Exemplo 2 -- Casca Cilíndrica}
  Raio \tabto{30mm}\emph{$r(x) = x$}

  \pause
  \vspace{\baselineskip}
  Circunferência \tabto{30mm}\emph{$2\pi x$}

  \pause
  \vspace{\baselineskip}
  Altura \tabto{30mm}\emph{$h(x) = x - x^2$}
\end{frame}

\begin{frame}{Exemplo 2 -- Volume}
  \begin{align*}
    \onslide<+->{V & = \int_0^1 2\pi x\left(x - x^2\right)\, dx                       \\[2mm]}
    \onslide<+->{  & = 2\pi \int_0^2 x^2 - x^3 dx                                     \\[2mm]}
    \onslide<+->{  & = 2\pi \left( \frac{x^3}{3} - \frac{x^4}{4}\right) \evalat{0}{1} \\[2mm]}
    \onslide<+->{  & = 2\pi \left( \frac{1}{3} - \frac{1}{4}\right) }
    \onslide<+->{  \;=\; 2\pi \frac{1}{12}}
    \onslide<+->{  \;=\; \frac{\pi}{6}}
  \end{align*}
\end{frame}

%------------------------------------------------------------------------------%
\section{Exemplo 3}

\begin{frame}{Exemplo 3}
  Encontre o volume do sólido obtido pela rotação da região
  delimitada por
  \[
    y = x - x^2
    \qquad\text{e}\qquad
    y = 0
  \]
  em torno da reta $x = 2$
\end{frame}

\framedgraphic{Exemplo 3 -- Esboço}{Volume_cascas_exemplo_3}{}

\begin{frame}{Exemplo 3 -- Casca Cilíndrica}
  Raio \tabto{30mm}\emph{$r(x) = 2-x$}

  \pause
  \vspace{\baselineskip}
  Circunferência \tabto{30mm}\emph{$2\pi (2-x)$}

  \pause
  \vspace{\baselineskip}
  Altura \tabto{30mm}\emph{$h(x) = x-x^2$}
\end{frame}

%------------------------------------------------------------------------------%
\begin{frame}{Exemplo 3 }
  \begin{align*}
    \onslide<+->{V & = \int_0^1 2\pi(2- x)\left(x - x^2 \right) dx         \\[2mm]}
    \onslide<+->{  & = 2\pi \int_0^1 x^3-3x^2 +2x \, dx                    \\[2mm]}
    \onslide<+->{  & = 2\pi \left( \frac{x^4}{4}- x^3+x^2 \right) \evalat{0}{1} \\[2mm]}
    \onslide<+->{  & = 2\pi \left( \frac{1}{4} - 1 + 1 \right) }
    \onslide<+->{  \;=\; \frac{\pi}{2}}
  \end{align*}
\end{frame}

%------------------------------------------------------------------------------%
\section{Exemplo 4}

\begin{frame}{Exemplo 4}
  Use cascas cilíndricas para encontrar o volume do sólido obtido pela rotação,\\
  em torno do eixo $x$ da região sob a curva
  \[y= \sqrt{x}\]
  de 0 até 1
\end{frame}

%\framedgraphic{Exemplo 4 -- Esboço}{Volume_cascas_exemplo_3}{}

\begin{frame}{Exemplo 4 -- Casca Cilíndrica}
  Raio \tabto{30mm}\emph{$r(y) = y$}

  \pause
  \vspace{\baselineskip}
  Circunferência \tabto{30mm}\emph{$2\pi y$}

  \pause
  \vspace{\baselineskip}
  Altura \tabto{30mm}\emph{$h(y) = 1-y^2$}
\end{frame}



%------------------------------------------------------------------------------%
\begin{frame}{Exemplo 4}
  \begin{align*}
    \onslide<+->{V & = \int_0^1 2\pi y(1-y^2) dy                                     \\[2mm]}
    \onslide<+->{  & = 2\pi \int_0^1 y-y^3 dy                                        \\[2mm]}
    \onslide<+->{  & = 2\pi \left(\frac{y^2}{2} - \frac{y^4}{4}\right) \evalat{0}{1} \\[2mm]}
    \onslide<+->{  & = 2\pi \left(\frac{1}{2} - \frac{1}{4} \right)                  \\[2mm]}
    \onslide<+->{  & = \frac{\pi}{2}}
  \end{align*}
\end{frame}

%------------------------------------------------------------------------------%
\section{Exemplo 5}

\begin{frame}{Exemplo 5}
  Encontre o volume do sólido de revolução gerado pela rotação da região\\
  delimitada pelas curvas
  \[
    x = (y-3)^2
    \qquad\text{e}\qquad
    x = 4
  \]
  em torno da reta $y=1$
\end{frame}

%\framedgraphic{Exemplo 4 -- Esboço}{Volume_cascas_exemplo_3}{}

\begin{frame}{Exemplo 5 -- Casca Cilíndrica}
  Raio \tabto{30mm}\emph{$r(y) = y-1$}

  \pause
  \vspace{\baselineskip}
  Circunferência \tabto{30mm}\emph{$2\pi (y-1)$}

  \pause
  \vspace{\baselineskip}
  Altura \tabto{30mm}\emph{$h(y) = 4-(y-3)^2$}
\end{frame}

%------------------------------------------------------------------------------%
\begin{frame}{Exemplo 5}
  \begin{align*}
    \onslide<+->{V & = \int_1^5 2\pi(y-1)\big(4-(y-3)^2\big) dy      \\[2mm]}
    \onslide<+->{  & = 2\pi \int_1^5 (y-1) \left(-y^2+6y-5\right) dy \\[2mm]}
    \onslide<+->{  & = 2\pi \int_1^5 -y^3+7y^2-11y+5 \, dy                \\[2mm]}
    \onslide<+->{  & = 2\pi \left( -    \frac{y^4}{4}
                                   + 7  \frac{y^3}{3}
                                   - 11 \frac{y^2}{2}
                                   + 5y \right) \evalat{1}{5}  \\[2mm]}
    \onslide<+->{  & = 2\pi \Bigg[\left( -\frac{5^4}{4}
                                   + 7    \frac{5^3}{3}
                                   - 11   \frac{5^2}{2}
                                   + 25   \right)
                                   - \left( -\frac{1}{4}
                                   +         \frac{7}{3}
                                   -         \frac{11}{2}
                                   + 5       \right) \Bigg] }
    \onslide<+->{  \;=\; \frac{128\pi}{3}}
  \end{align*}
\end{frame}

%------------------------------------------------------------------------------%
\section{Exemplo 6}

\begin{frame}{Exemplo 6}
  Use os dois métodos estudados para encontrar o volume do sólido
  gerado pela rotação da região delimitada pelas curvas
  \[
    y = 5-4x
    \qquad
    y = \sqrt{x}
    \qquad
    x = 0
  \]
  em torno do eixo $y$
\end{frame}

%\framedgraphic{Exemplo 6 -- Esboço}{Volume_cascas_exemplo_3}{}

%\begin{frame}{Exemplo 6 -- Casca Cilíndrica}
%  Raio \;\emph{$2-x$}
%
%  \pause
%  \vspace{\baselineskip}
%  Circunferência \;\emph{$2\pi (2-x)$}
%
%  \pause
%  \vspace{\baselineskip}
%  Altura \;\emph{$y = x-x^2$}
%\end{frame}
%
%\begin{frame}{Exemplo 6}
%  \begin{align*}
%    \onslide<+->{V & = \int_0^1 2\pi(2- x)\left(x - x^2 \right) dx         \\[2mm]}
%  \end{align*}
%\end{frame}



%------------------------------------------------------------------------------%
\begin{frame}
  \tableofcontents
\end{frame}

%------------------------------------------------------------------------------%
\end{document}
%------------------------------------------------------------------------------%
